Para el cálculo de integrales en regiones con simetría circular, se utiliza el sistema de coordenadas polares. Un punto $(x, y)$ se relaciona con $(r, \theta)$ a trav\'ez de  la transformaci\'on:

\begin{equation} \label{def:coord_polares}
\mathbf{T}:D^*=[0,R]\times[0, 2\pi] \subseteq\Rn{2} \to D\: : \: \mathbf{T}(r, \theta ) =( r \cos \theta , r \sin \theta ).
\end{equation} 

\begin{tarea} Probar que la transformaci\'on $T$ definida en (\ref{def:coord_polares}) est\'a bajo las las hi\'otesis del Teorema   \ref{ teo:cambio_variables_r2}.
\end{tarea}

\vspace{0.2cm}

Para el c\'alculo del \emph{Determinante Jacobiano} de $T$  primero  calculamos su  matriz diferencial $D_\mathbf{T}$, siendo
       \begin{equation*}
       \boldsymbol{D}_{\mathbf{T}}(r, \theta) 
       =
       \begin{pmatrix}
       \cos\theta & -r\sin\theta \\
       \sin\theta & r\cos\theta
       \end{pmatrix}
      \end{equation*}
   luego, el jacobiano es
       \begin{equation*}
       J_{\mathbf{T}}(r, \theta) = \det(\boldsymbol{D}_{\mathbf{T}})(r, \theta) = 
      r(\cos^2\theta + \sin^2\theta) = r,
      \end{equation*} 
  
por \'ultimo,  en nuestro caso,  como $r \geq 0$, tenemos
\begin{equation}\label{ejemplo_jacobiano_polares}
|J_{\mathbf{T}}| = r
\end{equation}

\clearpage

\textbf{Visualización del sistema de coordenadas polares como en esféricas y cilindricas}
\tikzexternaldisable % Desactiva la externalización para esta figura
\begin{figure}[H]
    \centering
    \begin{tikzpicture}

    %%% IZQUIERDA: caja D* en espacio (r, θ)
    \begin{axis}[
        name=izq,
        view={110}{25},
        width=7cm,
        height=7cm,
        xmin=-0.2, xmax=3.5,
        ymin=-0.2, ymax=7.5,
        zmin=-0.2, zmax=0.5,
        xlabel={$r$},
        ylabel={$\theta$},
        axis lines=center,
        ticks=none,
        every axis x label/.style={at={(ticklabel* cs:1.05)}, anchor=north east},
        every axis y label/.style={at={(ticklabel* cs:1.05)}, anchor=north west},
        every axis z label/.style={at={(ticklabel* cs:1.05)}, anchor=south},
    ]

    % Cara D* (θ ∈ [0, 2π], r ∈ [0, R])
    \addplot3[fill=magenta!30, opacity=0.8, draw=magenta!70, thick]
        coordinates {(0,0,0) (3,0,0) (3,6.28,0) (0,6.28,0) (0,0,0)};

    % Etiqueta D*
    \node[color=magenta!80!black, font=\Large]
        at (axis cs: 1.5, 3.14, 0) {$D^*$};

    % Cotas
    \node[below, color=black!60, font=\small]
        at (axis cs: 3, 0, 0) {$R$};
    \node[above right, color=black!60, font=\small]
        at (axis cs: 0, 6.28, 0) {$2\pi$};
    \node[left, color=black!60, font=\small]
        at (axis cs: 0, 0, 0) {$0$};

    \end{axis}

    % Label inferior izquierdo
    \node[anchor=north, font=\small] at (izq.south)
        {Espacio $(r,\theta)$};

    %%% FLECHA T
    \draw[-{Stealth[length=9pt, width=6pt]}, thick]
        ($(izq.east)+(0.4,0)$)
        -- node[above, font=\large]{$T$}
        ++(1.5,0);

    %%% DERECHA: disco D en espacio (x, y)
    \begin{axis}[
        name=der,
        at={($(izq.east)+(2.4cm,0)$)},
        anchor=west,
        view={110}{25},
        width=7cm,
        height=7cm,
        xmin=-2.8, xmax=2.8,
        ymin=-2.8, ymax=2.8,
        zmin=-0.2, zmax=0.5,
        xlabel={$x$},
        ylabel={$y$},
        axis lines=center,
        ticks=none,
        every axis x label/.style={at={(ticklabel* cs:1.05)}, anchor=north east},
        every axis y label/.style={at={(ticklabel* cs:1.05)}, anchor=north west},
        every axis z label/.style={at={(ticklabel* cs:1.05)}, anchor=south},
    ]

    % Disco D
    \addplot3[
        surf,
        domain=0:360,
        domain y=0:2,
        samples=30,
        samples y=5,
        fill=green!30,
        draw=green!55,
        fill opacity=0.75,
        draw opacity=0.9,
        variable=\t,
        variable y=\u
    ]
    ({\u*cos(\t)}, {\u*sin(\t)}, {0});

    % Borde del disco (r = R)
    \addplot3[
        draw=green!70,
        thick,
        domain=0:360,
        samples=60,
        variable=\t
    ]
    ({2*cos(\t)}, {2*sin(\t)}, {0});

    % Etiqueta D
    \node[color=green!50!black, font=\Large]
        at (axis cs: 0.4, 0.4, 0) {$D$};

    \end{axis}

    % Label inferior derecho
    \node[anchor=north, font=\small] at (der.south)
        {Espacio $(x,y)$};

    \end{tikzpicture}

    \caption{Visualización del sistema de coordenadas polares. 
    El bloque $D^*$ a la izquierda representa la región en el espacio 
    $(r,\theta)$ con sus límites típicos, mientras que el disco 
    $D$ a la derecha muestra cómo se transforma esa región en el 
    espacio $(x,y)$ mediante la transformación $T$.}
    \label{fig:polares_cambio_figura}
\end{figure}
\tikzexternalenable % Vuelve a activar la externalización para las demás figuras


\begin{example}
Hallar el área de la región
\[
D = \{(x,y) \in \mathbb{R}^2 : 1 \leq x^2 + y^2 \leq 4\}.
\]
\end{example}

\textbf{Soluci\'on:} Primero, dibujemos el \'area pedida 


\tikzexternaldisable
\begin{center}
\begin{tikzpicture}[scale=0.9, >=stealth]

  %%% IZQUIERDA: Dominio D* en espacio (r, theta)
  % Relleno del rectángulo D*
  \fill[blue!20] (1, 0) rectangle (2, 6.28);
  % Borde del rectángulo D*
  \draw[blue!70!black, thick] (1, 0) rectangle (2, 6.28);
  
  % Ejes r y theta
  \draw[->] (-0.5, 0) -- (3.0, 0) node[right] {$r$};
  \draw[->] (0, -0.5) -- (0, 7.0) node[above] {$\theta$};
  
  % Marcas y etiquetas en eje r
  \draw (1, 0.08) -- (1, -0.08) node[below] {\small $1$};
  \draw (2, 0.08) -- (2, -0.08) node[below] {\small $2$};
  
  % Marcas y etiquetas en eje theta
  \draw (0.08, 0) -- (-0.08, 0) node[left] {\small $0$};
  \draw (0.08, 6.28) -- (-0.08, 6.28) node[left] {\small $2\pi$};
  
  % Etiqueta D*
  \node[color=blue!80!black, font=\large] at (1.5, 3.14) {$D^*$};
  \node[below=0.6cm, font=\small] at (1.25, 0) {Espacio $(r,\theta)$};

  %%% FLECHA DE TRANSFORMACIÓN T
  \draw[->, ultra thick] (3.6, 3.14) -- node[above, font=\large] {$T$} (5.1, 3.14);

  %%% DERECHA: Imagen D (Corona Circular) en espacio (x, y)
  \begin{scope}[shift={(8.3, 3.14)}]
    % Relleno de la corona D
    \begin{scope}
      \clip (-2.6, -2.6) rectangle (2.6, 2.6);
      \fill[blue!20] (0, 0) circle (2);
      \fill[white]   (0, 0) circle (1);
    \end{scope}
    
    % Círculos bordes
    \draw[blue!70!black, thick] (0, 0) circle (2);
    \draw[blue!70!black, thick] (0, 0) circle (1);
    
    % Ejes cartesiano x e y
    \draw[->] (-2.6, 0) -- (2.7, 0) node[right] {$x$};
    \draw[->] (0, -2.6) -- (0, 2.7) node[above] {$y$};
    
    % Marcas en los ejes
    \foreach \x in {-2,-1,1,2}
      \draw (\x, 0.06) -- (\x, -0.06) node[below] {\small $\x$};
    \foreach \y in {-2,-1,1,2}
      \draw (0.06, \y) -- (-0.06, \y) node[left] {\small $\y$};
      
    % Etiqueta D
    \node[color=blue!80!black, font=\large] at (1.55, -0.4) {$D$};
    \node[below=0.6cm, font=\small] at (0, -2.1) {Espacio $(x,y)$};
  \end{scope}

\end{tikzpicture}
\end{center}
\tikzexternalenable

La región $D$ es una corona circular de radio interior $1$ y radio exterior $2$, luego podemos usar el cambio de coordenadas polares definidos en (\ref{def:coord_polares}) con la restricci\'on de dominio correspondiente, es decir, 
 \begin{equation*}
\mathbf{T}:D^*=[1,2] \times[0, 2\pi] \subseteq\Rn{2} \to D\: : \: \mathbf{T}(r, \theta ) =( r \cos \theta , r \sin \theta ).
\end{equation*} 

Siguiendo la Definición \ref{def:area}, el área se calcula como:
\begin{equation*}
A(D) = \iint_D 1\, dA
\end{equation*}

luego, por el Teorema de cambio de variable  (\ref{ teo:cambio_variables_r2}) y recordando (\ref{ejemplo_jacobiano_polares}) 
$$ \iint_D 1\, dA = \iint_{D^*} r\, dr\, d\theta= \int_0^{2\pi} \int_1^2 r\, dr\, d\theta $$ 

 por \'ultimo, usando  el Teorema de Fubbinni  (\ref{col:fubini}) 

\begin{equation*}
\int_0^{2\pi} \int_1^2 r\, dr\, d\theta = \int_0^{2\pi} \Big(  \int_1^2 r\, dr\,  \Big) d\theta  =  \int_0^{2\pi}  \Big( \frac{r^2}{2}\Big|_1^2 \Big)  d\theta  = \int_0^{2\pi}  \frac{3}{2}  d\theta =3\pi.
\end{equation*}

\textbf{Cambio de coordenadas Lineales} \textcolor{red}{Agregar}

